% THETA3-A7 contribution by /root/theta_032, 21 November 2026.
% GPT-6 Astra Pro; inherited reasoning-effort label not separately exposed.

\subsection{THETA3-A7: an independent wedge is a genuine single theta lift}\label{A7}

\paragraph{Setup and statement.}
Let $V$ have rank $n+1$, signature $(n,1)$ at one real place and
positive definite signature elsewhere, with $[F:\Q]\geq2$.
Let $P=P_1\perp\cdots\perp P_n$ be the skew-Hermitian multiplicity
space of the first $n$ types in A6. Its active component is definite.
Put
\[
 G=\U(V),\qquad S=\U(P),\qquad T=\prod_{i=1}^n\U(P_i).
\]
The quotient $[S]=S(F)\backslash S(\A_F)$ is compact: a nonzero
global isotropic vector would remain isotropic in the definite
active component. This assertion does not require $S$ to be compact
at every real place.

Fix compatible theta splittings, a finite Schwartz input, and the
degree-$n$ holomorphic Fock cocycle. At the active place this is
the determinant coefficient of the wedge of $n$ degree-one
cocycles; at the other places take the specified Gaussian
coefficients. Let
\[
 \mathcal K_h=\theta_V(\,\cdot\,,h;\Phi)\in\mathcal E,\qquad
 \mathcal E=H^0(X,\Omega_X^n),
\]
in the equivariant adelic realization on a fixed sufficiently small
compact Shimura level, including its finitely many components.
The notation includes the indicated differential-form projection.
Fix a product of unitary source characters
$\pi(t)=\prod_i\pi_i(t_i)$ on $[T]$, and use the
van Hoften--Sempliner convention
\[
 F=\int_{[T]}\mathcal K_t\,\pi(t)\,dt.
\]
For a decomposable input the tensor rule and degree-$n$ alternation
identify $F$ with the wedge of the corresponding $n$ integrated
degree-one lifts. For a finite sum of inputs it is the corresponding
finite sum of wedges.

There is a smooth automorphic function $v_\pi$ on $[S]$, of finite
level and finite maximal-compact type, such that the exact identity
\begin{equation}\label{A7-lift}
 F=\int_{[S]}\theta_V(\,\cdot\,,h;\Phi)\,
                                  \overline{v_\pi(h)}\,dh
\end{equation}
holds with the displayed measures. Thus the independent wedge is
an actual single theta lift from $\U(P)$. The representative
$v_\pi$ is allowed to depend on the source characters and on the
kernel. No character independence is asserted.

\paragraph{Why the kernel has finite rank.}
First we verify that $\mathcal K_h$ belongs to the same
finite-dimensional holomorphic space for every $h\in S(\A_F)$.
At infinity the source Weil action commutes with the target
$G$-action, its maximal compact action, and its relative Lie
algebra differential. Applying the source action to a closed
cochain of Hodge type $(n,0)$ therefore gives another closed
cochain of that type. The corresponding differential form is
$\bar\partial$-closed and hence holomorphic. The original cochain
is the wedge of the degree-one holomorphic cocycles; its
closedness follows from the derivation rule for the cochain
differential. The explicit source is BMM, Lemma 5.4 and
Proposition 5.17(1),(4), pp.43 and 48. No stable-range assertion
about generation in degree $n$ is used.

At finite places the source and target Weil actions commute.
Every element of a compact open target group fixing the original
finite input consequently fixes its source translate. The target
level is thus independent of $h$. The chosen target splitting
and canonical compact type are retained. Compactness of $X$
implies that $\mathcal E$ is finite-dimensional.

The theta series and its derivatives converge uniformly on compact
sets of representatives, by Schwartz decay and the positive
Gaussian estimates. This gives a smooth $\mathcal E$-valued
function of $h$. It is left $S(F)$-invariant by rational theta
automorphy. A sufficiently small finite source subgroup fixes
$\Phi_f$, giving a common finite level. At infinity a polynomial
times a Gaussian is finite under a source maximal compact
subgroup, so the resulting coefficient functions have finite
maximal-compact type. These statements use compact representative
sets for the anisotropic adelic quotients, not a finite-dimensional
assertion about all functions with one compact type on a
noncompact real group.

Choose a Petersson orthonormal basis
$f_1,\ldots,f_D$ of $\mathcal E$, with inner products conjugate
linear in the first variable. Then
\begin{equation}\label{A7-kernel}
 \mathcal K_h=\sum_{j=1}^D f_j\,e_j(h),
 \qquad e_j(h)=\langle f_j,\mathcal K_h\rangle_G .
\end{equation}
The functions $e_j$ are smooth automorphic functions on $[S]$
with the level and compact-type properties just proved.
Let $\mathcal C=\operatorname{span}_{\C}\{e_1,\ldots,e_D\}$.
It is finite-dimensional. Its $L^2([S])$ inner product is positive
definite: a nonzero smooth function is nonzero on an open set
of positive Haar measure. In particular the passage from smooth
functions to their $L^2$ classes is injective here.

\paragraph{The exact Riesz representative.}
Define a complex-linear functional on $\mathcal C$ by
\[
 \Lambda_\pi(e)=\int_{[T]}e(t)\pi(t)\,dt.
\]
The inclusion $T\hookrightarrow S$ gives a well-defined continuous
map $[T]\to[S]$. The torus quotient is compact, so this integral
exists for every $e\in\mathcal C$. Its boundedness as a functional
on $\mathcal C$ follows from finite dimensionality; no assertion
about boundedness of pointwise restriction on an unrestricted
$L^2$ space is needed.

By finite-dimensional Riesz representation there is a unique
$v_\pi\in\mathcal C$ such that
\[
 \Lambda_\pi(e)=\langle v_\pi,e\rangle_S
               =\int_{[S]}\overline{v_\pi(h)}e(h)\,dh
 \quad(e\in\mathcal C).
\]
Substitute \eqref{A7-kernel} and interchange the finite sum and
the integrals:
\[
 F=\sum_jf_j\Lambda_\pi(e_j)
   =\int_{[S]}\left(\sum_jf_je_j(h)\right)
                                  \overline{v_\pi(h)}\,dh.
\]
This is \eqref{A7-lift}. If $F\ne0$, then
$\Lambda_\pi$ is nonzero on $\mathcal C$, so $v_\pi\ne0$.
For an orthonormal basis $\eta_1,\ldots,\eta_r$ of $\mathcal C$,
the representative and its norm are explicitly
\[
 v_\pi=\sum_{a=1}^r
              \overline{\Lambda_\pi(\eta_a)}\,\eta_a,
 \qquad
 \|v_\pi\|_S^2=\sum_{a=1}^r|\Lambda_\pi(\eta_a)|^2.
\]
These are analytic finite-dimensional quantities; their values
or normalizations have not been proved algebraic.
\hfill$\square$

The assertion is an equality of theta integrals with a smooth
automorphic input. It does not select a prescribed irreducible
source representation. When $P$ is totally definite, the entire
archimedean source group is compact; the stated finite source
level and finite archimedean type then place $v_\pi$ in a
finite-dimensional automorphic module, which has a finite
spectral decomposition. For the more general anisotropic source,
finite rank of the particular coefficient space $\mathcal C$
is what proves the identity.

\paragraph{The signature condition for a totally definite middle source.}
At a compact real place $\tau$, the sign of the Hermitian form on
$P_i$ is
$-d_i(\tau)\operatorname{sgn}(\operatorname{Im}\tau(\delta))$.
Consequently $P$ is definite there if and only if all
$d_i(\tau)$ agree. It is already definite at the active place.
Thus the actual multiplicity space is totally definite if and
only if
\[
 v_1=v_2=\cdots=v_n
\]
as complete type functions. For $n=2$, a place where $v_1$
and $v_2$ differ gives signature $(1,1)$ on $P$. Multiplying
the form by a scalar can reverse the signs but cannot make that
component definite.

There are always nonzero wedges from one repeated type: by A6's
fibre-spanning proof, one theta family supplies $n$ forms with
independent values. Such a wedge has a totally definite middle
source in \eqref{A7-lift}. But for the particular balanced tensor
of A6, equality of the first $n$ types forces equality of all
$2n$ types. Indeed at each $g$, the first sum is either $n$ or
$-n$; a sum of $n$ numbers in $\{1,-1\}$ attains that value
only when each summand has the same sign. Hence every last-list
$v_i$ equals the common first-list function. The selected
balanced tensor line is then a product of divisor lines.

This separates an arbitrary nonzero canonical wedge from the
particular wedge needed to reach a selected non-Lefschetz
Weil component. Changing to isogenous CM factors does not
remove this condition: the signs follow the complete types of
the selected eigenlines, including their embedding labels.

In the direct $m=2$ instantiation of the recovered
\texttt{MainX}, the middle group is required to be totally
definite. Formula \eqref{A7-lift} therefore provides the actual
middle source only in the just-described definite case.
For a general exceptional balanced input, it supplies a genuine
single lift from an anisotropic group with compact quotient,
but not the required totally definite middle space. It does
not rule out another source realization of the same canonical
automorphic representation. A second reverse lift, its chosen
vector, and the complete \texttt{LawsX} and arithmetic link data
remain separate requirements for that iterated comparison.
